\subsection{The finite connection formulas}\label{sec:connections}
Here $P$ and $Q$ name only the finite paired boundary fields, all source
fields being in $\dom(J_*)$, so $J_*(P)=Q$. The source endpoint $h$
is not one of those fields. The side record $S(P)$ may extend beyond
$L_h$; no assertion $J_*(S(P))=S(Q)$ is needed. Its closed fields and
certificates are in $L_c$. A side field can lie between two boundary
item fields; the required inequalities are
\begin{equation}\label{eq:7-2}
 c(X)<\text{variable side field}<c(Y)\quad\text{in }R(X,Y),
\end{equation}
together with the actual complete field order. Fixed earlier ports
are not variable side fields. Open old items retain their old labels,
Gen, actual graphs, exact domains and work-level relations; their new
private heights are not yet asserted to be complete.

\begin{lemma}[One-supplier finite connection]\label{lem:connection}
Suppose the prepared source supplies true $R(P,Q)$ and
$\operatorname{Init}(P)$, correct low-side certificates in $L_c$, and
no omitted external requests. For any finite number of copies and any
$\beta<a$, the original $J_*$ alone supplies an ordered chain of low
boundaries, all its required side records and incoming edges, and the
final low certificates. All new fields are below $a$; the final closed
certificates belong to $L_a$.
\end{lemma}
\begin{proof}
The high witness $X=P$ gives in $L_v$
\begin{equation}\label{eq:7-3}
 \exists X\,[\beta<c(X)<c\ \wedge\operatorname{Init}(X)
                                      \wedge R(X,Q)].           
\end{equation}
Reflect through $J_*$, whose fixed external parameters lie below $a$.
It gives $A_0$ below $a$, with Init$(A_0)$ and $R(A_0,P)$. Equation
(\ref{eq:7-2}) puts all its side fields below $a$, even if they interleave
the item fields of $A_0$.

Suppose the last low boundary is $A$ and $R(A,P)$ has already been
supplied. Its source certificate lies in $L_h$, hence also in $L_v$,
where its bounded check remains correct. Use $Y=P$ as the high witness
for
\begin{equation}\label{eq:7-4}
 \exists Y\,[c(A)<c(Y)<c\ \wedge R(A,Y)\wedge R(Y,Q)].       
\end{equation}
Keep the previously required low fields and earlier connections in
the formula. They lie below $a$ and are fixed by $J_*$. Certificates
not yet below $a$ remain existential fields; they are not falsely
treated as fixed parameters. Reflection gives $A<B<a$ with $R(A,B)$
and $R(B,P)$. Retain all needed relations $R(A_i,P)$, as well as the
consecutive connections. Their real records preserve each open item's
stability to $L_a$. Repeat this finite step as many times as the
specified native expansion requires.

Finally use $Z=P$ and the actual $S(P)$ to cap the entire low chain.
In the high existence statement require every final closed certificate
separately to belong to $L_c$. This statement has a true witness:
the earlier connection certificates are in $L_h\subset L_c$, and the
original closed side certificates are in $L_c$ by hypothesis.
Since $J_*(L_a)=L_c$, reflection gives the final cap below $a$ and
all the required certificates in $L_a$. Retained incoming graphs are
part of the resulting data; none has been used as an additional
stage-1 supplier.

Choose the common work model with the same $J_*$, by reflecting the
joint statement (\ref{eq:5-2}), which includes both its strength and the actual
open-item elementary relations. Then discard $P,Q,h,v,J_*$ and the
high records. Future backward travel uses only the actual low
connections and their Init source. The induction has taken finitely
many downward uses of one original map and has not used the output
of a previous native step as an extra assumption.
\end{proof}

The preparation supplies the hypotheses of this lemma, not merely an
isomorphism of displayed heights. Proper rootOnly gaps are filled by
whole images along actual old routes, and complete both packages are
kept. Old coverage blocks remain as side records. Every direct incoming
edge at a right boundary belongs to $R$, including targetOnly terminal
edges with their exact domains. Lower-side closed graphs have targets
below $c$, so the tight certificate theorem places their data in $L_c$.
Right-side graphs have targets below $v$, so their data are in the
actual $L_v$. A certificate for a relation to the high boundary $P$
may remain in $L_h$ until the cap. These are different layer bounds;
placing every intermediate certificate in $L_c$ without this argument
would lose the tight bare-endpoint case.

Stage 0, when permitted, is a finite source-prefix synchronization of
the original last request. Freeze that request before adding anything.
At its first missing proper pair, the actual consumer route reaches
a complete package strictly before the original last column. If the
root is an ancestor, the strict-before-containing-head conclusion of
Theorem~\ref{thm:invariants} supplies this bound. Transport that entire
old package; subsequent transports use its continuous whole images,
each still before the same original last column. The next unpaid
original proper-root gap moves forward. No internal modification is
performed before transporting its already changed enclosing package.
When the frozen request has no rootOnly, none of these stage-0
operations is used; its existing qualified coverage and both skeleton
provide the appropriate records directly.
